\documentclass[a4paper]{article}
\usepackage{graphicx}
\usepackage{amsmath,amssymb}
\usepackage{hyperref}

\title{Autobiographic}
\begin{document}
    \maketitle
    According to \cite{brewer-1986-what-is-autobiographical-memory}, defined as memory for information related to the self additionally he includes the first three of the following as different forms of autobiographical memory
    \begin{enumerate}
        \item \textit{Personal memory}. I might experience a mental image that corresponds to a particular episode in my life, such as the time in California, while visiting Mt. Palomar, when I made a snowball and threw it down the trail at my sons, starting a snowball fight.
        \item \textit{Autobiographical fact}. I might simply recall the fact (with no accompanying imagery) that I drove to Mt. Palomar on another occasion when I was a senior in high school.
        \item \textit{Generic personal memory}. I might have a mental image of what, in general, it was like to drive north on Highway 1 in California’s Big Sur region: an image that does not appear to be of any specific moment during the drive but is a generic view from the driver’s seat of the car with the sharp jagged cliffs to my right and the Pacific stretching out to the horizon on my left.
        \item \textit{Semantic memory}. I might recall (with no accompanying imagery) that in California in 1978 the voters passed a limitation on property tax called Proposition 13.
        \item \textit{Generic perceptual memory}. I might have a visual image of the outline of the shape of the state of California.
    \end{enumerate}
    
    
    
    
    \url{https://en.wikipedia.org/wiki/Autobiographical_memory} says:
    Autobiographical memory (AM) is a memory system consisting of episodes recollected from an individual's life, based on a combination of episodic (personal experiences and specific objects, people and events experienced at particular time and place) and semantic (general knowledge and facts about the world) memory. It is thus a type of explicit memory.
    
    The autobiographical knowledge base contains knowledge of the self, used to provide information on what the self is, what the self was, and what the self can be. This information is categorized into three broad areas:
    \begin{itemize}
        \item lifetime periods
        \item general events
        \item event-specific knowledge
    \end{itemize}
    
    \section{Types}
        \begin{itemize}
            \item \textbf{Biographical or Personal:}
            These autobiographical memories contain personal biographical information, such as where a person was born or the names of their parents.
        
            \item \textbf{Copies vs.\ Reconstructions}
            \begin{itemize}
                \item \textbf{Copies:} Vivid autobiographical memories of an experience with a considerable amount of visual and sensory-perceptual detail. These memories may vary in authenticity.
                \item \textbf{Reconstructions:} Autobiographical memories that are not direct reflections of raw experiences, but are rebuilt by incorporating new information or interpretations formed in hindsight.
            \end{itemize}
        
            \item \textbf{Specific vs.\ Generic}
            \begin{itemize}
                \item \textbf{Specific:} Memories that contain detailed recollection of a particular event (event-specific knowledge).
                \item \textbf{Generic:} Memories that are vague and contain little detail beyond the type of event that occurred. Episodic autobiographical memories may also fall into this category when one remembered event represents a series of similar events.
            \end{itemize}
        
            \item \textbf{Field vs.\ Observer}
            \begin{itemize}
                \item \textbf{Field memories:} Memories recalled from the original perspective, experienced from a first-person point of view.
                \item \textbf{Observer memories:} Memories recalled from a perspective outside oneself, as if observing the event from a third-person viewpoint. Older memories are often recalled in this form and are more frequently reconstructions, while field memories tend to be more vivid and similar to copies.
            \end{itemize}
        
            \item \url{https://en.wikipedia.org/wiki/Autobiographical_memory#Types}
        \end{itemize}
    
     \bibliographystyle{plain}
     \bibliography{/home/donkarlo/Dropbox/repo/nd_shared_research/bibliography/bibliography.bib}
\end{document}